\section{原则一：选择正确的工具集}

更多的工具\textbf{并不总是}带来更好的结果。这一原则是文章所有原则中最具战略性的一条。

\subsection{Agent 的认知特性与传统软件的差异}

Agent 拥有与传统软件不同的 \textbf{affordance}（感知到的可操作性）。LLM Agent 的"上下文"是有限且昂贵的，而传统计算机内存则廉价且充足。

文章用\textbf{通讯录搜索}来类比说明这一差异：

\begin{table}[H]
\centering
\caption{通讯录搜索：传统软件 vs Agent}
\begin{tabular}{p{0.15\textwidth}p{0.35\textwidth}p{0.35\textwidth}}
\toprule
 & \textbf{传统软件} & \textbf{LLM Agent} \\
\midrule
方法 & 遍历联系人列表，逐个检查 & 必须将信息加载到有限的上下文窗口中 \\
\midrule
成本 & 内存廉价，遍历高效 & 每个 Token 都消耗上下文空间 \\
\midrule
最佳策略 & \texttt{list\_contacts} 后过滤 & \texttt{search\_contacts} 直接定位 \\
\bottomrule
\end{tabular}
\end{table}

\begin{warningbox}{常见反模式：简单包装已有 API}
最常见的错误是将已有的软件功能或 API 端点简单包装为 Tool，而不考虑这些工具是否适合 Agent 使用。这相当于强迫 Agent 用"逐页阅读通讯录"的方式来查找联系人——即暴力搜索（brute-force search），而自然的方式是直接按字母跳转到相关页面。
\end{warningbox}

\subsection{工具整合的策略}

工具可以整合功能，在底层处理多个离散操作或 API 调用。文章给出了三组对比：

\begin{table}[H]
\centering
\caption{工具设计：碎片化 vs 整合}
\begin{tabular}{p{0.42\textwidth}p{0.48\textwidth}}
\toprule
\textbf{碎片化设计（不推荐）} & \textbf{整合设计（推荐）} \\
\midrule
\texttt{list\_users} + \texttt{list\_events} + \texttt{create\_event} & \texttt{schedule\_event}：自动查找可用时间并安排事件 \\
\midrule
\texttt{read\_logs} & \texttt{search\_logs}：只返回相关日志行及上下文 \\
\midrule
\texttt{get\_customer\_by\_id} + \texttt{list\_transactions} + \texttt{list\_notes} & \texttt{get\_customer\_context}：一次性编译客户的所有相关信息 \\
\bottomrule
\end{tabular}
\end{table}

每个工具应有\textbf{清晰、独特的用途}。工具应该让 Agent 以类似人类的方式分解和解决任务，同时减少中间输出对上下文的消耗。

过多工具或功能重叠的工具会分散 Agent 的注意力，使其偏离高效策略。精心、有选择地规划工具集可以带来显著的回报。

\subsection{本章小结}

选择正确的工具集是一个战略性决策。核心原则是：(1) 不要简单包装已有 API；(2) 优先实现面向高影响力工作流的少量精心设计的工具；(3) 通过工具整合减少 Agent 的认知负担和上下文消耗；(4) 确保每个工具有清晰独特的用途。


